% ============================================================
% PARTE IX - METODO DI LAVORO (seguito)
% CAPITOLO 44 - COME AFFRONTARE IL BASSO DATO
% ============================================================
\chapter{Come affrontare il basso dato (numerato e non numerato)}
\chaptermark{Come affrontare il basso dato}
\label{chap:metodo_lavoro}

L'armonizzazione di un basso dato è una prova tipica degli esami di armonia. Il capitolo dà l'ordine in cui applicare le regole già studiate, in sette passi che valgono per il basso numerato e per quello senza numeri. Ogni passo rimanda ai capitoli dove la teoria è spiegata. La scelta degli accordi sul basso senza numeri è nel Capitolo \ref{chap:scegliere_numerica}; la Sezione \ref{sec:bassi_svolti} applica il metodo a quattro bassi svolti.

Nel basso numerato (o figurato) la numerica dice già quali intervalli costruire sopra ogni nota, e il lavoro consiste nel tradurre i numeri in accordi con una buona condotta delle voci. Nel basso senza numeri le informazioni si ricavano dal basso stesso, con i passi 2--4 e con i criteri del Capitolo \ref{chap:scegliere_numerica}. Gli altri passi sono uguali nei due casi.

\section{Leggere la numerica}
\label{sec:lettura_numerica_metodo}

Ogni cifra è un intervallo contato dal basso reale, non dalla fondamentale dell'accordo, usando le note dell'armatura (Osservazione \ref{oss:numerica_relativa}). Le alterazioni scritte nella numerica prevalgono sulle note della scala. Un $\sharp$ o un $+$ da solo alza la terza dell'accordo.

\textit{Teoria:} Sezione \ref{subsec:cifratura_accordi} (triadi), Sezione \ref{sec:lettura_numerica} (settime e alterazioni).

\section{Analisi preliminare}
\label{sec:fase1_numerato}

Prima di scrivere accordi si legge l'intero basso in silenzio e ci si chiede:
\begin{itemize}
	\item qual è l'armatura, e quali sono la prima e l'ultima nota (di norma la tonica);
	\item se il modo è maggiore o minore (la sensibile alzata ricorre regolarmente nel minore armonico, e nel minore melodico salgono anche il VI grado);
	\item qual è il metro, e quindi quanti accordi si prevedono (di norma uno per accento principale);
	\item se ci sono alterazioni non previste dall'armatura: possono indicare la sensibile del minore, una dominante secondaria, un prestito dal modo parallelo o una modulazione (Sezione \ref{sec:alterazioni_modulazione});
	\item se alcune note sono estranee (passaggio, volta, pedale) e quindi non portano accordi propri;
	\item se ci sono note lunghe o legate, che possono portare più accordi (Sezioni \ref{sec:basso_legato} e \ref{sec:note_lunghe}).
\end{itemize}

\textit{Teoria:} Capitoli \ref{chap:tonalita}, \ref{chap:tonicizzazione}, \ref{chap:modulazione}, \ref{chap:noteestranee}.

\section{Segmentazione fraseologica}
\label{sec:fase2_numerato}

Si divide il basso in unità di senso compiuto e si pianificano i punti di riposo.
\begin{itemize}
	\item Si cercano le cadenze: una nota lunga dopo un tratto in movimento, V--I con salto di quarta o quinta, il basso che si ferma sulla dominante. La numerica, quando c'è, le rivela: \texttt{4 3} indica un ritardo cadenzale, \texttt{7} sulla dominante la settima in cadenza.
	\item In un basso di 8 misure si pianifica la cadenza di metà (sospesa su V, oppure imperfetta) e la cadenza finale (perfetta). In un basso di 4 misure vale lo stesso criterio.
	\item Se la numerica non specifica la risoluzione della dominante si sceglie la forma completa o incompleta secondo il contesto.
\end{itemize}

\textit{Teoria:} Capitoli \ref{chap:cadenze} e \ref{chap:struttura_frase}, Definizione \ref{def:v7_completa_incompleta}; Tabella \ref{tab:forme_cadenzali} per le forme cadenzali secondo la durata del basso.

\section{Classificare ogni nota del basso}
\label{sec:fase3_numerato}

Si passa nota per nota, assegnando funzione e stato dell'accordo. Per ogni nota ci si chiede, in ordine:
\begin{enumerate}
	\item Qual è il grado nella scala?
	\item Cade su un tempo forte o debole? I tempi forti portano di norma stati fondamentali.
	\item Come si muove rispetto alla nota precedente e alla successiva: grado congiunto (Regola dell'Ottava), salto (stato fondamentale), nota ribattuta o tenuta?
	\item Fa parte di un tratto riconoscibile: cadenza, progressione, scala?
	\item Un'alterazione indica una dominante secondaria, un prestito o una modulazione?
	\item Quale accordo, fra quelli possibili, è più coerente con il contesto? Una nota del basso da sola non indica nessuna funzione: la decidono le note vicine.
\end{enumerate}

Per ogni nota i candidati sono quattro (fondamentale, terza, quinta o settima di un accordo) e si scelgono con i filtri del Capitolo \ref{chap:scegliere_numerica}, Regola \ref{reg:filtri_scelta}.

Per i tratti scalari lunghi o con note ripetute si segue la regola seguente.

\begin{regola}[Strategia per l'armonizzazione del basso]
	\label{reg:strategia-armonizzazione}
	\begin{itemize}
		\item Movimenti scalari brevi (2--4 note): Regola dell'Ottava.
		\item Movimenti scalari lunghi (5 o più note): si valutano le progressioni.
		\item Salti: prevalentemente stati fondamentali.
		\item Note ribattute o legate: stesso accordo, oppure settima al basso se la nota scende di grado.
		\item Quando si riconosce una progressione, si applica la stessa armonizzazione a tutte le ripetizioni del modello.
	\end{itemize}
\end{regola}

\textit{Teoria:} Capitolo \ref{chap:scegliere_numerica}, Capitoli \ref{chap:funzioni_armoniche}--\ref{chap:ottava_preliminare}, \ref{chap:progressioni_armoniche}--\ref{chap:moduli}, \ref{chap:regola_ottava_completa}, \ref{chap:tonicizzazione}.

\section{Costruire le voci superiori}
\label{sec:fase4_numerato}

Stabilita l'ossatura armonica si scrive a quattro parti.

\begin{regola}[Priorità nella scrittura delle voci superiori]
	\label{reg:priorita_scrittura_numerato}
	\begin{enumerate}
		\item Vettori obbligati: la sensibile sale alla tonica, la settima scende di grado. Hanno priorità assoluta.
		\item Note comuni: si tengono ferme nella stessa voce.
		\item Grado congiunto: le altre voci si muovono per la distanza minima.
		\item Moto contrario: si preferisce fra basso e soprano, per evitare quinte e ottave nascoste.
	\end{enumerate}
\end{regola}

Si scelgono i raddoppi (di norma la fondamentale; mai sensibile o settima) e la posizione del soprano, cercando una linea cantabile. Si controlla che le voci restino nella zona medio-alta dell'estensione e ben separate fra loro.

\textit{Teoria:} Capitoli \ref{chap:condotta_voci}--\ref{chap:errori_condotta}.

\section{Gestire le dissonanze}
\label{sec:fase5_numerato}

Quando la numerica o il piano armonico prevedono una settima, un ritardo o una nota estranea:

\begin{regola}[Trattamento delle dissonanze nella realizzazione]
	\label{reg:dissonanze_realizzazione}
	\begin{enumerate}
		\item Si individua la nota dissonante, cioè quella che crea la dissonanza rispetto al basso.
		\item Si verifica la preparazione: la nota era consonante nell'accordo precedente, nella stessa voce? È obbligatoria per \sigla{II}{7}{}, \sigla{IV}{7}{} e per i ritardi, non per la \sigla{V}{7}{}. Se la dissonanza è la nota del basso, la preparazione è il basso legato.
		\item Si pianifica la risoluzione: la nota scende di grado. Qual è la sua destinazione nell'accordo successivo?
		\item Si controlla che le risoluzioni simultanee non producano quinte o ottave parallele.
	\end{enumerate}
\end{regola}

\textit{Teoria:} Capitoli \ref{chap:noteestranee}--\ref{chap:ritardi}, \ref{chap:settima_dominante}--\ref{chap:ritardi_complessi}.

\section{Revisione finale}
\label{sec:fase6_numerato}

Si applica la lista di controllo (Regola \ref{reg:checklist_errori}, Sezione \ref{sec:checklist}), rileggendo ogni coppia di voci isolata (S+B, A+T, S+A, T+B) prima di valutare l'insieme. Un buon elaborato oltre a non contenere errori ha un discorso musicale coerente: la numerica vincola la scelta degli accordi, mentre la condotta delle voci e la qualità melodica del soprano restano libere.

\section{Quattro bassi svolti}
\label{sec:bassi_svolti}

I quattro esercizi seguenti sono bassi senza numeri. Per ciascuno si legge la sequenza delle note, si fissano le cadenze e si sceglie un accordo per nota con i criteri del Capitolo \ref{chap:scegliere_numerica}. Sotto le note del basso, nelle figure, stanno il grado e la numerica.

\subsection*{Primo basso: Do maggiore con 6/4 cadenzale}

Basso: Do--La--Si--Do $\mid$ Fa--Sol--Sol--Do.

La tonalità è Do maggiore e il basso ha due semifrasi di quattro note. La prima torna al Do e prolunga la tonica; la seconda chiude con il salto di quinta Sol--Do, quindi con una cadenza perfetta.

Do--La--Si--Do sale per grado da La a Do, e la Regola dell'Ottava assegna \sigla{IV}{6}{} al La e \sigla{V}{6}{5} al Si: la formula è \sigla{I}{}{}--\sigla{IV}{6}{}--\sigla{V}{6}{5}--\sigla{I}{}{} (Sezione \ref{sec:formule_frase}). Nella seconda semifrase il Fa è un \sigla{IV}{}{} dopo la tonica, e il Sol ribattuto prima del Do è \sigla{I}{6}{4}--V.

Il soprano canta Mi--Fa--Fa--Mi e poi Fa--Mi--Re--Do. Il Fa del \sigla{V}{6}{5} è la settima e scende al Mi. Nel \sigla{I}{6}{4} il Do e il Mi scendono a Si e Re, mentre il Sol resta fermo.

\begin{esempio}[Basso Do--La--Si--Do--Fa--Sol--Sol--Do]
	\label{ex:basso_svolto_1}
	\musica{es_cap44_01}
\end{esempio}

\subsection*{Secondo basso: La minore con settima al basso}

Basso: La--La--Sol$\sharp$--La $\mid$ Re--Mi--La.

La minore, due semifrasi. Nella prima il La è legato e il basso scende poi al Sol$\sharp$, che è la sensibile: il La può essere la settima di un \sigla{iiø}{2}{}, già preparata dalla tonica precedente, e scende proprio al Sol$\sharp$. Il Sol$\sharp$ al basso è la terza di V, quindi \sigla{V}{6}{}. La prima semifrase è \sigla{i}{}{}--\sigla{iiø}{2}{}--\sigla{V}{6}{}--\sigla{i}{}{}.

La seconda semifrase ha Re--Mi--La, la cadenza S-D-T. Il Re può essere \sigla{iv}{}{} o \sigla{ii°}{6}{}; si sceglie \sigla{iv}{}{} perché il basso sale di grado alla dominante e il passaggio da sottodominante a dominante è diretto.

Il soprano canta Mi--Fa--Mi--Mi e poi La--Sol$\sharp$--La. Il Fa della prima semifrase è la quinta diminuita del \sigla{iiø}{2}{} e scende al Mi del \sigla{V}{6}{}; il Sol$\sharp$ del soprano nella seconda semifrase sale al La.

\begin{esempio}[Basso La--La--Sol$\sharp$--La--Re--Mi--La]
	\label{ex:basso_svolto_2}
	\musica{es_cap44_02}
\end{esempio}

\subsection*{Terzo basso: modulazione da Do a Sol}

Basso: Do--La--Fa$\sharp$--Sol $\mid$ Do--Re--Sol.

Il basso comincia in Do maggiore e contiene un Fa$\sharp$, che in Do non è nella scala. Il Fa$\sharp$ è la terza di Re--Fa$\sharp$--La--Do, quindi \sigla{V}{6}{5} di Sol, e il basso sale subito al Sol. Dopo il Sol il basso continua con Do--Re--Sol, la cadenza \sigla{IV}{}{}--V--\sigla{I}{}{} in Sol maggiore: la nuova tonalità è confermata da una cadenza, quindi è una modulazione (Regola \ref{reg:ton_mod_basso}).

Il pivot è il \sigla{vi}{}{} di Do maggiore, che in Sol maggiore è il \sigla{ii}{}{}. La prima parte è \sigla{I}{}{}--\sigla{vi}{}{}--\sigla{V}{6}{5}--\sigla{I}{}{}, la seconda \sigla{IV}{}{}--\sigla{V}{7}{}--\sigla{I}{}{}.

Il soprano tiene il Do nei primi tre accordi (tonica, \sigla{vi}{}{}, \sigla{V}{6}{5}). Nel \sigla{V}{6}{5} il Do è la settima e scende al Si della nuova tonica; la cadenza finale ha Do--La--Sol.

\begin{esempio}[Basso Do--La--Fa$\sharp$--Sol--Do--Re--Sol: modulazione Do M $\rightarrow$ Sol M]
	\label{ex:basso_svolto_3}
	\musica{es_cap44_03}
\end{esempio}

\subsection*{Quarto basso: sottodominante minore e note lunghe}

Basso: Do--Fa--Fa--Sol--Sol--Do.

Do maggiore, con due note ripetute (Fa e Sol) prima della tonica finale. Il Sol ripetuto è il segnale del 6/4 cadenzale: \sigla{I}{6}{4}--V (Sezione \ref{sec:note_lunghe}). Sul Fa ripetuto si possono scegliere due accordi: \sigla{IV}{}{} e poi \sigla{iv}{}{}, preso a prestito dal modo minore (Sezione \ref{sec:iv_maggiore}). La successione è \sigla{I}{}{}--\sigla{IV}{}{}--\sigla{iv}{}{}--\sigla{I}{6}{4}--V--\sigla{I}{}{}.

Il soprano canta Mi--Fa--Fa--Mi--Re--Do. Il La del contralto nel \sigla{IV}{}{} diventa La$\flat$ nel \sigla{iv}{}{} e scende al Sol del \sigla{I}{6}{4}. Nel 6/4 il Mi scende al Re e il Do scende al Si, mentre il Sol resta fermo.

\begin{esempio}[Basso Do--Fa--Fa--Sol--Sol--Do, con sottodominante minore]
	\label{ex:basso_svolto_4}
	\musica{es_cap44_04}
\end{esempio}
